\head{本轮核查：局部超 Tor、正则序列与超平面截面}
\noindent\textbf{核查依据与界限。}本章直接对照上传的 \texttt{EGAIII.2.pdf}、\texttt{EGAIV.1.pdf}、\texttt{EGAVaiello.pdf} 的连续正文及页面图像。新续译的范围是 III（6.4.1）--（6.6.8）、$0_{\mathrm{IV}}$（15.1.5）--（15.2.5），以及 V 预备笔记第一篇§§3--6。下面的反例、补证和修订建议由本次译校撰写，不是原著正文。V 的核查只针对所上传的英文译编本，未见对应法文手写原稿，不能据此分配原作者、英译者和排版者的责任。

本章有三项针对所写断言的明确反例（B2-E01--E03）；其中B2-E03取决于上传英文的字面读法，已说明可能缺少的等价连接词。另有两项证明步骤问题（B2-G01--G02），以及排印、定义约定与未完成稿件的注意事项。\textbf{这些编号不应相加称为发现了同样数量的错误定理。}原报告中的条目保留在后续两编，作为历史核查记录，并非本次重新逐一认证。

\egasection{本轮定位表}
\begin{longtable}{@{}p{.15\textwidth}p{.27\textwidth}p{.50\textwidth}@{}}
\toprule 编号 & 底本条目／PDF页 & 被检查的具体问题\\\midrule\endhead
B2-E01 & IV.1，15.1.15／16--17 & 任意基变换保持正则性的条件漏列最终商模平坦。\\
B2-E02 & V，4.4／15 & 用闭集本身的维数，不能推出一般态射的泛截面连通性。\\
B2-E03 & V，5.5／17 & 无嵌入相伴循环不使任意超平面自动成为非零因子。\\
B2-G01 & IV.1，15.1.9／15 & 对非忠实张量函子作单射性消去；另给完整补证。\\
B2-G02 & III.2，6.5.14(i)／19 & 假设没有保证底概形 $S$ 诺特；证明实际上不需要它。\\
B2-T01--T05 & III.2、IV.1 & “双射／单射”、漏非单位、Čech指标基数、三个／四个及引文问题。\\
B2-T06--T10 & V，§§3--5 & 参数空间、载体、局部平凡化对象及等号错置。\\
B2-H01--H09 & 见正文脚注及本章 & 历史术语、定义与继承假设，不作为无条件反例。\\
B2-P01--P04 & V，3.1、3.1.1、4.4、5.6 & 非约化情形、缺证明、可疑英译、未明确的引文。\\
\bottomrule
\end{longtable}

\audit{B2-E01}{任意基变换不保持所有正则序列}{原假设不足；具体反例}
\source{\texttt{EGAIV.1.pdf}，PDF第16页，原书第17页，（15.1.15）；证明续于下一页}
\textbf{原文实际条件。}设 $A\to B$ 为交换代数，$M$ 为 $B$-模，$f_1,\ldots,f_n$ 为 $M$-正则序列。底本要求
\[
 M\Big/\Bigl(\sum_{j=1}^{i-1}f_jM\Bigr)
 \quad(1\leq i\leq n)
\]
均为平坦 $A$-模，随后声称对任意 $A\to A'$，变换后的序列仍正则。注意求和上限确实是 $i-1$；原页图像已核对，不是从OCR下标猜出的结论。它只要求 $M_0,\ldots,M_{n-1}$ 平坦，没有要求 $M_n$ 平坦，其中
\[
 M_i=M/(f_1,\ldots,f_i)M,
 \qquad M_0=M.
\]

\textbf{反例及逐项核验。}取任意域 $k$，令
\[
 A=B=k[t],\quad M=A,\quad n=1,\quad f_1=t,
 \quad A'=A/(t)=k.
\]
$A\to B$ 为恒等同态，$M$ 是自由 $A$-模，因而满足原文唯一要求的平坦性；又因 $k[t]$ 为整环，乘以 $t$ 的映射 $M\to M$ 单射，所以 $(t)$ 为 $M$-正则序列。基变换后
\[
 B'=k,
 \qquad M'=k\ne0,
 \qquad f'_1=t\otimes1=0.
\]
乘以 $f'_1$ 是 $k\to k$ 的零映射，其核为整个 $k$，故不单射。全部所写假设成立，结论失败。

\textbf{出错的精确步骤。}原证明一开始使用
\[
 0\longrightarrow M\xrightarrow{f_1}M\longrightarrow M/f_1M\longrightarrow0
\]
在张量 $A'$ 后仍正合，其理由是 $M/f_1M$ 平坦。然而在 $n=1$ 的原条件中，这个商模并没有被要求平坦。上述例子张量后的左侧变成 $k\xrightarrow{0}k$，正好显出所缺的一项。

\textbf{一个明确成立的修订版本。}保留正则序列假设，要求每个 $M_i$（$0\leq i\leq n$）都 $A$-平坦。则对任意 $A\to A'$，由 $\Tor_1^A(M_i,A')=0$，各列
\[
 0\to M_{i-1}\otimes_AA'
 \xrightarrow{f_i\otimes1}M_{i-1}\otimes_AA'
 \to M_i\otimes_AA'\to0
\]
都正合。而
\[
 M_{i-1}\otimes_AA'
 \simeq M'/(f'_1,\ldots,f'_{i-1})M',
\]
故变换后的序列正则。各 $M_i\otimes_AA'$ 又是平坦 $A'$-模，所以同样保留全部商模的平坦性。这个版本足以修补原论证；这里不宣称它是所有情形中最弱的假设。

\textbf{不应误伤下一命题。}（15.1.16）的证明从最终商模平坦，利用局部平坦性准则逐步倒推所有 $M_i$ 平坦；因此它的实际应用满足上述修订条件。本反例不能单独用来否定（15.1.16）。现代的相对正则浸入理论同样明确保留被切出的子概形对底空间的平坦性；参见 Stacks Project，Tag 063P，Lemma 31.23.1及其后定义。这个外部对照不是反例的来源，反例本身完全由上面的计算验证。

\audit{B2-G01}{不能从非忠实张量后的单射性消去张量}{证明步骤缺乏依据；给出不依赖该步骤的补证}
\source{\texttt{EGAIV.1.pdf}，PDF第15页，原书第16页，（15.1.9）}
\textbf{原文推理。}它将典范同态 $\varphi$ 分解后，得到某个 $\varphi'\otimes1$ 为同构，继而说 $\varphi'$ 必单射。这里的张量因子含 $A/(f_n)$，它一般不忠实；单靠“张量后单射”不能作这一推论。

例如 $A=k[t]$，满同态 $q:A\to A/(t)$ 不单射，而
\[
 q\otimes_AA/(t):k\longrightarrow k
\]
是恒等同构。本例只反驳这一普遍的消去规则；它不是（15.1.9）整个命题的反例，因为后者另有分离性等假设，必须在证明中实际使用这些假设。

\subsec{在原命题假设下的完整替代论证}
设 $J=(f_1,\ldots,f_n)$，且
\[
 \varphi:(M/JM)[T_1,\ldots,T_n]\xrightarrow{\sim}\gr_J(M)
\]
是定义拟正则性的典范同构。假设 $M$ 及原命题列出的各前缀商模都对 $J$-预进制拓扑分离。证明 $(f_1,\ldots,f_n)$ 为 $M$-正则序列。

\textbf{第一步：$f_1$ 对 $M$ 单射。}
在 $\gr_J(M)$ 上，由乘以 $f_1$ 诱导的次数为1的映射，在 $\varphi$ 下就是多项式变量 $T_1$ 的乘法。无论系数模是否有挠，乘以多项式变量总是单射：它把每个单项式系数移到一个不同的单项式位置。

若 $0\ne x\in M$，分离性 $\bigcap_rJ^rM=0$ 保证存在最大的非负整数 $r$，使 $x\in J^rM$；也就是
$x\in J^rM\setminus J^{r+1}M$。其初始类 $\bar x\in\gr_J^r(M)$ 非零，故 $T_1\bar x\ne0$，即
\[
 f_1x\notin J^{r+2}M.
\]
特别地 $f_1x\ne0$，所以 $f_1$ 为 $M$-正则元素。

\textbf{第二步：计算相交。}对每个 $r\geq1$，有
\begin{equation}\label{eq:new-intersection}
 f_1M\cap J^rM=f_1J^{r-1}M.
\end{equation}
从右到左的包含由 $f_1\in J$ 直接得到。反过来，若 $f_1y\in J^rM$ 而 $y\notin J^{r-1}M$，则 $y\ne0$，可取其有限 $J$-阶 $s\leq r-2$。第一步关于初始类的计算给出 $f_1y\notin J^{s+2}M$，却由 $r\geq s+2$ 得到 $f_1y\in J^rM\subseteq J^{s+2}M$，矛盾。因此 $y\in J^{r-1}M$，证明了\eqref{eq:new-intersection}。

\textbf{第三步：模去第一项。}令 $N=M/f_1M$，给它取商 $J$-进制滤过。对 $r\geq1$，自然满同态
$J^rM/J^{r+1}M\to J^rN/J^{r+1}N$ 的核为
\[
 \frac{J^{r+1}M+(J^rM\cap f_1M)}{J^{r+1}M}
 =\frac{J^{r+1}M+f_1J^{r-1}M}{J^{r+1}M},
\]
即 $T_1\gr_J^{r-1}(M)$。第零次核为零，因为 $f_1M\subseteq JM$。因此
\[
 \gr_J(N)\simeq\gr_J(M)/T_1\gr_J(M)
 \simeq(M/JM)[T_2,\ldots,T_n].
\]
同时 $N/JN\simeq M/JM$，所以右端正是剩余序列 $f_2,\ldots,f_n$ 对 $N$ 的拟正则性典范同构。

\textbf{第四步：归纳。}原假设保证 $N$ 以及继续模去 $f_2,\ldots$ 后的前缀商模均分离。因此对 $N$ 重复前三步，得到 $f_2$ 在 $N$ 上单射，再得到 $f_3$ 在下一个商模上单射，依次进行到 $f_n$。这正是原文的 $M$-正则序列定义。$n=1$ 时第一步已经给出全部结论。逆向“正则推出拟正则”仍使用前面的（15.1.8），无需改动。

\textbf{结论。}命题可以成立，但原印刷证明中“因此 $\varphi'$ 必单射”的句子不能仅靠非忠实的张量消去证明。上述补证显式使用分离性，避免了这个跳步；未要求 $M$ 本身完备，也未偷加诺特或有限生成条件。

\audit{B2-G02}{局部超 Tor 的证明不需要底空间诺特}{原证明附带断言错误；核心结论可保留}
\source{\texttt{EGAIII.2.pdf}，PDF第19页，原书第20页，（6.5.14）证明(i)}
（6.5.13）的条件(i)只要求 $X$ 和 $Z=X\times_SY$ 局部诺特，且 $X\to S$ 平坦。缩到仿射情形后，底本说“$S,X,Z$ 都诺特”；其中关于 $S$ 的断言没有依据。

取
\[
 A=k[x_1,x_2,x_3,\ldots],\quad K=\operatorname{Frac}(A),\quad
 S=\Spec A,\quad X=Y=\Spec K.
\]
$A$ 不诺特，因为理想链 $(x_1)\subsetneq(x_1,x_2)\subsetneq\cdots$ 不稳定。$K$ 是 $A$ 的局部化，因此 $X\to S$ 平坦。又
\[
 K\otimes_AK\simeq K,
\]
故 $X$、$Z=\Spec K$ 都诺特。条件(i)成立，$S$ 却不诺特。

\textbf{保留原结论的证明。}令仿射坐标环分别为 $A,B,C$，$D=B\otimes_AC$。$X$ 和 $Z$ 诺特表示 $B,D$ 诺特。有限 $B$-模 $F$ 有由有限自由 $B$-模组成的解消：先选有限生成元得到 $B^{r_0}\to F$，其核因 $B$ 诺特而有限，再重复这一过程。由于 $B$ 对 $A$ 平坦，这一解消的各项也对 $A$ 平坦。因此用它与有限 $C$-模 $G$ 作外张量即可计算 $\Tor^A_n(F,G)$。各项 $B^{r_i}\otimes_AG$ 为有限 $D$-模；因为 $D$ 诺特，其核、像和同调仍为有限 $D$-模。对应的层为凝聚层。这个论证没有使用 $A$ 诺特，正好说明多写的 $S$ 诺特可从证明中删去，不必加强原定理的假设。

\audit{B2-E02}{一般态射下按原空间维数控制断开集不够}{连通性断言有明确反例}
\source{\texttt{EGAVaiello.pdf}第15页，（4.4）；连同其后的原稿自我批评}
先把明显的参数记号整理为：对每个分支 $X_i$ 有 $\dim\overline{f(X_i)}\geq2$，参数在 $P^\vee$ 中。即便如此，原文对一般态射的 $(C_{m+1})\Rightarrow$一般截面满足$(C_m)$ 仍不成立。

\subsec{空间和态射}
取代数闭域 $k$，并置
\[
 X=\Spec k[u,v,w,z]/(uv),\qquad
 f:X\longrightarrow\mathbf P^2_k,
 \quad f=[1:u+v:(u+v)w].
\]
第一齐次坐标恒为1，所以这个态射处处有定义。$X$ 有两个不可约分支
\[
 X_1=V(u)\simeq\mathbf A^3_k,
 \qquad X_2=V(v)\simeq\mathbf A^3_k,
\]
交集 $D=V(u,v)\simeq\mathbf A^2_k$，由 $w,z$ 作坐标。在 $X_1$ 上，$f$ 的两个仿射坐标为 $(v,vw)$；对任意第一坐标非零的点 $(a,b)$，可取 $v=a,w=b/a$，故像稠密于 $\mathbf A^2$。在 $X_2$ 上同理。因此两个像闭包都为二维，满足原文各分支的像维数条件。

\subsec{\texorpdfstring{$X$ 满足 $(C_1)$}{X 满足 (C1)}}
设 $T\subset X$ 为维数至多1的闭集。因 $\dim X_i=3$，它不能包含任何一个 $X_i$，所以 $X_i\setminus T$ 为非空不可约开子集，因而连通。又因 $\dim D=2$，$T$ 不能包含 $D$，于是
\[
 (X_1\setminus T)\cap(X_2\setminus T)=D\setminus T\ne\varnothing.
\]
两个连通子空间有非空交，它们的并 $X\setminus T$ 连通。这验证了 $(C_1)$；因为 $k$ 已代数闭，无须另作代数闭包扩张。

\subsec{泛超平面逆像却不连通}
令 $K=k(a,b)$，以 $1+aT_1+bT_2=0$ 表示对偶射影平面一个标准开图上的泛超平面。其逆像的坐标环为
\[
 R=K[u,v,w,z]/\bigl(uv,\ 1+(u+v)(a+bw)\bigr).
\]
令 $s=u+v$、$c=a+bw$，则在 $R$ 中 $sc=-1$，故 $s,c$ 都可逆。由 $uv=0$ 可得
\[
 e_1=\frac{u}{s},\qquad e_2=\frac{v}{s},\qquad
 e_1+e_2=1,
 \quad e_1e_2=0,
 \quad e_i^2=e_i.
\]
两个幂等元都非零，因为分别置 $v=0$ 或 $u=0$ 后得到的商环均为
\[
 K[w,z,1/(a+bw)]\ne0.
\]
故 $\Spec R$ 分解为两个非空不交开闭部分；事实上
\[
 R\simeq K[w,z,1/(a+bw)]\times K[w,z,1/(a+bw)].
\]
所以 $Y_\eta$ 已不连通，更不可能满足几何 $(C_0)$。这里的断开由删除空集即可看出；也可直接用这两个非空开闭部分表述，避免空集维数的约定差异。

这还反驳了“好参数包含稠密开集”的表述：当超平面第一系数非零、第三系数非零时，同样的计算给出两个非空分支。这些参数形成一个稠密开集，因而不可能与另一个全部给出连通截面的稠密开集并存。

\textbf{缺失的几何控制。}$D$ 本身二维，却被 $f$ 压到单个点 $[1:0:0]$。泛超平面避开这个点，便把原先用来连接两个分支的全部交集去掉。原文用闭集自身的维数控制断开行为，不能排除这种高维闭集被压缩的情况。原稿随后建议改用像的维数，正是针对这一机制；本报告只验证上述反例，不宣称已经替它证明所有拟议的推广。

\textbf{英译问题的界限。}原稿后面英文写“not valid as such if $f$ is quasi-finite”，疑似缺少限制词；上面的 $f$ 并不拟有限，故不能用它判定拟有限版本真假，更不能静默将英文改成另一句话后说那就是原作者原话。

\audit{B2-E03}{无嵌入相伴循环不使每个超平面自动正则}{上传英文的字面断言不成立；可能遗漏等价连接词}
\source{\texttt{EGAVaiello.pdf}第17页，（5.5）}
这里的 $V$ 已在（5.4）定义为使截面 $\phi_\xi$ 对 $F_{k(\xi)}$ 正则的参数集合。英文（5.5）先假设 $F_s$ 无嵌入相伴循环，随后说“we see immediately that $\xi\in V$”。如果把它读成无条件的推出，便错误。

令 $S=\Spec k$，$X=P=\mathbf P^2_k$，$f=\id$，$L\subset P$ 为直线，$i:L\hookrightarrow P$ 为闭浸入，取 $F=i_*\OO_L$。该凝聚模的唯一相伴素循环是 $L$，没有嵌入分支。选取参数 $\xi$ 使 $H_\xi=L$。定义该超平面的截面在 $F$ 上作用为零，因此
\[
 F(-1)\xrightarrow{\phi_\xi}F
\]
为零映射，而 $F(-1)\ne0$；它不单射，所以 $\xi\notin V$。

\textbf{应当保留的判据。}在所述无嵌入相伴循环的有限型情形，$\xi\in V$ 应等价于：$\Supp F_{k(\xi)}$ 的每个不可约分支均\emph{不包含于} $H_\xi$。它不是从“无嵌入循环”自动推出的结论，而是还须检查的相交条件。该判据来自非零因子避开全部相伴素理想的刻画。英文的“which means also”可能原本想表达这种等价关系；本报告只纠正所上传载体的字面读法，不判定法文手稿的原句。

同段声称模去 $\phi$ 后 Cohen--Macaulay 性保留时，也必须先保证该截面正则。上述例子同时说明，不能把这一前提省略。

\head{本轮排印核对、条件提醒与待核事项}
\audit{B2-T01}{“双射”应为“单射”}{证明中的词语误植}
\source{IV.1，PDF第14页，原书第15页，（15.1.8）证明末段}
原文连续两次使用 bijective，认为末项 $f_n$ 的数乘在相关商分次模上双射；正则性只保证单射。取 $A=k[x,y]$、$M=A$、平凡滤过，$f_1=x,f_2=y$。在 $A/(x)=k[y]$ 上乘以 $y$ 单射但不满射，因为1不在像中。证明只需要单射，用自由模张量保持单射后应用（15.1.6）即可。这里不反驳命题结论。

\audit{B2-T02}{再次漏列“非单位”}{重复的例子条件遗漏}
\source{IV.1，PDF第16页，原书第17页，（15.1.12）（iii）}
该例沿用 $K/A$，写每个非零 $f$ 均不对 $\gr(M)$ 正则。若 $f=1$，数乘恒等，显然正则；此时所比较的两个分次模都为零，也失去反例作用。要求 $f\ne0$ 且不是单位后，原构造才工作。这是上批E02在新段落中的重复表现，没有重新计成独立的主要定理错误。

\audit{B2-T03}{Čech 上链的指标基数符号错了}{索引排印问题}
\source{III.2，PDF第21--22页，原书第22--23页，（6.6.4）}
对普通 Čech 上链 $C^r$，每个分量由 $r+1$ 个开集的交标记。原文使用同调方向的 $h\leq0$，相应上链为 $C^{-h}$，所以指标集应有 $1-h$ 个元素，而不是 $-(h+1)$ 个。仅代 $h=0$ 就能检查：$C^0$ 由单个开集上的截面组成，指标数是1；原公式给出$-1$，没有意义。

第二个方向同理是 $1-k$。若 $h_1+h_2=h$，两边指标总数为
\[
 (1-h_1)+(1-h_2)=2-h,
\]
不是原文随后印的 $-h-2$。这里没有改变复形的微分符号；只是让指标基数与原文的 $C^{-h}$ 约定相符。

\audit{B2-T04}{四个谱函子误写三个}{计数文字误植}
\source{III.2，PDF第23页，原书第24页，（6.6.7）}
同一句先列 $(a),(b),(a'),(b')$，随后写“这三个谱函子”。应为四个；原比较论证分别适用于所列四个，再用于 $(c),(d)$。不是谱序列收敛性的反例。

\audit{B2-T05}{Bourbaki引文内部不一致}{待逐版核对的引文问题}
\source{IV.1，PDF第13页，原书第14页，（15.1.6）}
其中一次引用为第III章§3第8小节，附近相关引用为§2第8小节。记录这一内部不一致，不以之断言数学命题错误。本次没有取得所引Bourbaki具体版次进行逐页校勘，因此不声称该书所有版本的页码、节次都已确认。

\audit{B2-T06}{不可约分支链的首项与并集记号}{V稿中的局部记号整理}
\source{V，第13页，（3.3）证明}
“$X_0$ and $X'$”应据链的两端读作 $X_0=X'$；后面混写的 $Y_\eta=X_\eta$，按同段定义应写 $Y_\eta=\bigcup_iX_{i\eta}$。译文将这两处按其同句定义明确排出，并逐处标注，未以这些明显的字母问题否定连通性证明。

\audit{B2-T07}{变动参数与层的载体必须分开}{V稿§4的对象类型和下标问题}
\source{V，第13--15页，§4开头及（4.2）--（4.4）}
$F$ 在 $X$ 上，$G$ 在 $Y$ 上；变动纤维应写 $Y_\xi,G_\xi,k(\xi)$，不能在随 $\xi$ 变化的集合定义中一直写泛点 $\eta$。超平面的参数属于 $P^\vee$，不属于 $P$。在“对每个不可约分支”的条件中，像的维数应逐个对 $f(X_i)$ 检查。这些类型整理本身并不修复B2-E02的实质性反例。

\audit{B2-T08}{平凡化的是可逆模而非任意模}{V稿§5的载体、基域与同态记号问题}
\source{V，第16页，（5.2）}
一般模 $F$ 不必局部同构于 $\OO_X$，例如秩2的自由模就不能；可以局部平凡化的是可逆模 $L$。截面 $\phi:\OO_X\to L$ 对偶地给出 $L^{-1}\to\OO_X$，然后张量 $F$ 得到 $F\otimes L^{-1}\to F$。因此底本缺失的另一张量因子是这个对偶同态，而不是预设一个全局处处存在的倒数函数。

同页从 $k(s)$ 扩张到 $k(\xi)$ 后，载体应为 $X_{k(\xi)}$，并先把原层限制到 $X_s$。不能在等式右端仍写 $X_{k(s)}$ 而忽略该扩张。

\audit{B2-T09}{纤维参数的基点写错}{V稿（5.3）的类型问题}
\source{V，第16页，（5.3）}
要取 $P^\vee\to S$ 在 $s$ 上的纤维，应有 $s\in S$，不是底本的 $s\in P^\vee$。所讨论的开集属于 $X_{P^\vee}$；$G$ 若仍作为 $Y$ 上的层，应限制到 $Y\cap U$，或者明确先作闭浸入直接像。它们不改变正则性与平坦性判据，但使用时必须说清对象。

\audit{B2-T10}{“相等”与下一句结论相反}{V稿证明中的等号误植}
\source{V，第17页，（5.3）（b）证明}
原文写 $Z_{i,k(\eta)}\cap Y=Z_{i,k(\eta)}$，却据此说截面不在任何相伴循环上恒零。正确所需的是不等号：若整个循环包含在 $Y$ 中，截面反而在其上恒零。原稿所指的维数理由也只支持真交，而不支持整条循环被包含。译文保留原等号并作注，没有静默修改。

\egasection{条件与定义：B2-H01--H09}
\textbf{B2-H01。}（III，6.5.5）的“上同调维数”是各茎环作为环的全局同调维数上界，不是任意层在拓扑空间上的上同调消失阶。不能见到相同术语就混用两类维数。

\textbf{B2-H02。}（III，6.5.8）中的仿射平坦性应经环同态 $A\to B$ 翻译：拟凝聚层 $\widetilde M$ 对 $S=\Spec A$ 平坦，当且仅当 $M$ 是平坦 $A$-模。底点局部化 $M\otimes_AA_s$ 与任取的上方点茎 $M_x$ 不是同一个记号。

\textbf{B2-H03。}本书（15.1.7）不要求最终商模非零，所以（15.1.12）中 $(1,0)$ 可以正则，$(0,1)$ 却不正则。其他书可能另加最终商非零的条件，不能据其他定义把原文例子判错。Stacks Project，Tag0AUH，Definition10.68.1采用额外非零条件，并紧接着明确提示有的作者不采用它；本次只作定义差异对照。

\textbf{B2-H04。}（15.1.17）使用Nakayama时，宜明确对有限 $B$-模 $R/\sum_if_iM$ 操作；剩余域张量条件使这个商模被 $\mm_AB$ 生成，而局部同态保证 $\mm_AB\subseteq\mm_B$，故商模为零。正则性则来自（15.1.16），不是单靠Nakayama推出。

\textbf{B2-H05。}（15.1.20）的“suite de composition”在此上下文表示有限滤过，商项同构于指定的 $M/(f_1,\ldots,f_n)M$；它们不必是简单模。按现代某些书的“合成列”狭义定义理解会平添错误条件。

\textbf{B2-H06。}（15.2.3）证明的是在 $V(J)$ 外拟正则性自动成立。不要据其“研究兴趣”的末句误读成有序正则性也自动成立；$(0,1)$ 就反驳后一个读法。

\textbf{B2-H07。}V（3.1.1）脚注猜测“primary”可能是代数闭。实际上传的 EGA IV 第二部分（4.3.1），PDF第55页，定义为较大域中所含的最大\emph{可分代数}扩张等于底域。取特征 $p>0$ 下 $k=\mathbf F_p(t)$、$L=k(t^{1/p})$，便见准素不要求相对代数闭。这一定义核对已完成，但IV第二部分正文没有因此被计为已经连续翻译。

\textbf{B2-H08。}V（4.3）印为“特征0并且非分歧”，前面（3.2）为“特征0或者非分歧”。两个充分条件同时加上会使假设更强，不单凭这种差别判定结论错误，也不静默改写连接词。

\textbf{B2-H09。}V（6.1）括号中的 $X$ 对 $k$ 固有不能删除。取
\[
 X=\Spec\mathbf C[x,y,z]/(xy-1)\simeq\mathbf G_m\times\mathbf A^1
 \hookrightarrow\mathbf A^3\hookrightarrow\mathbf P^3.
\]
它不可约二维，有限闭集不能将它断开；浸入的有限像闭集只能是有限点集，所以满足其余有关连通性和像维数的条件。然而超平面 $x+y=0$ 的截面由
$x^2+1=0$ 给出，分解为 $x=\sqrt{-1}$ 与 $x=-\sqrt{-1}$ 两条不交仿射直线。它不连通。这里 $X$ 不固有，故这只说明条件不可丢，不是对保留括号假设的（6.1）的反例。

\egasection{仍需保留的缺文与未核实事项：B2-P01--P04}
\textbf{B2-P01。}V（3.1）只假设几何不可约，直接称 $X$ 有“函数域”时，需要先说明传到 $X_{\mathrm{red}}$，或用泛点剩余域。该修补适用于拓扑不可约性论证，不是增强为几何整。

\textbf{B2-P02。}V（3.1.1）在本稿中没有证明；作者明确把选择证明出处的工作留给编稿者，英译者只列出几项参考。译文保留这一事实，没有把空缺改写成“证毕”。外部的Bertini--Zariski理论与本载体是否自足是两个问题。

\textbf{B2-P03。}V（4.4）的英文自我批评中“if $f$ is quasi-finite”与上下文可疑，尚未核对法文手稿；AQT的含义及末尾拉丁语短句也不强行猜定。B2-E02只建立一个可直接验证的一般态射反例，不借此代替来源学考证。

\textbf{B2-P04。}V（5.6）引用的12.1.1、12.1.4未在本节明确列出性质清单，英译编者自己说12.1.4核对不上。原证明还从涉及余核 $G$ 的结论转向 $F_{P^\vee}|_U$，须逐项核对所引命题。本次保留其开放问题状态，未凭通常的纤维性质开放性补写一个新的通用定理。

\egasection{本轮结论与引用方式}
三项明确反例中，B2-E01针对正则序列基变换命题的漏项，B2-E02针对一般态射的连通性断言，B2-E03针对英译稿中可疑连接词的字面读法。两项G类只否定所指证明步骤或附带断言，不能改说成定理反例。T类、H类、P类分别用于排印、条件和未完成事项；原稿的“参考待补”并不自动等于相应结论为假。

\textbf{本轮外部核验来源。}直接底本为本章开头列出的三份上传PDF；准素扩张的定义另核对上传 \texttt{EGAIV.2.pdf}（4.3.1）。外部仅作限定交叉核验：Stacks Project，Tag0AUH（Regular sequences）和Tag063P（Relative regular immersions），访问于2026年11月14日。可按稳定Tag在该项目检索：
\begin{center}
\url{https://stacks.math.columbia.edu/tag/0AUH}\\
\url{https://stacks.math.columbia.edu/tag/063P}
\end{center}
本章的具体反例和替代证明在上文完整写出，不以网页的权威性代替论证。
\par
